\chapter{Introduction}
\label{ch:introduction}

\section{Introduction}
\label{sec:introduction}

Communication is crucial for all people regardless of their ages. Therefore, the development of speech at early ages is an important factor affecting children's ability to communicate with other people and learn various information. According to statistics, approximately 7\% of children in the age range of 3--17 have been diagnosed with some type of disorder regarding their voice, speech, or language in the last year \cite{nidcd_statistics}. When children are experiencing speech difficulties, they may be unable to express themselves properly because of lowered self-esteem and the absence of contacts with peers. It would be helpful for parents to take their children to see a speech therapist; however, the effectiveness of the work with children greatly depends on the practice that children make outside therapy hours. For instance, Allen (2013) compared two groups of preschool children participating in phonological intervention. Children were divided into two categories according to the number of their practice sessions: three times a week and one time per week. Those who practised three times showed great results concerning the percentage of consonants correct while those who trained only once made similar results as the control group \cite{allen2013intervention}.

When put into the real world however, this creates a resourcing and availability challenge: therapy sessions are limited in time and frequency and at-home practice is often difficult to sustain consistently or perform correctly. Children simply lack motivation to do it or lose it due to the monotonous aspect of the training process. Even parents or caretakers may lack the understanding, resources, and structure necessary for proper home practice of their child. The discrepancy between treatment sessions and home practice is a practical issue in software engineering: whether such a software solution is capable of addressing it.

The bachelor project focuses on exploring the design and development of a software system which consists of two components aimed at helping children with speech disabilities through practice. This involves developing a mobile application through which the children can complete their speech-related tasks as well as a web-based therapist's dashboard, through which the speech therapists would be able to assign tasks and monitor the progress of their patient. The exercise content in the mobile application is inspired by established literature on articulation disorders \cite{todorova_articulation}, and includes diverse ways of practicing.

The proposed system is not intended to replace professional speech therapists but rather aims to act as a supplementary tool that makes structured speech practice more accessible between therapy sessions. The project not only looks at implementing this system but also takes into consideration the overall software engineering needed to implement this system successfully.

The remainder of this report is structured as follows: \Cref{ch:project-background} presents the project background and defines the scope of the work. \Cref{ch:problem-statement} introduces the problem statement and connects each problem to the proposed system features. \Cref{ch:implementation} presents the implementation, including the technologies, system components, application flow, exercise system, and results. \Cref{ch:software-engineering-practices} discusses the software engineering practices used throughout the project. \Cref{ch:discussion} reflects on the effectiveness of the solution, challenges encountered, technical debt, and lessons learned. Finally, \Cref{ch:conclusion} concludes the report by summarising the project and its main outcomes.
